%%=======================================================================
% !Mode:: "TeX:UTF-8"
% !TEX program  = PdfLaTeX
%%=======================================================================
\documentclass{beamer}
\usetheme{thubeamer}
\graphicspath{{figures/}}
\include{macros}

\usepackage{listings}
\usepackage{booktabs}
\usepackage{tikz}
\usetikzlibrary{shapes.geometric, arrows.meta, positioning, fit, calc}

% Define Rust language for listings
\lstdefinelanguage{Rust}{
  keywords={fn,let,mut,if,else,match,return,loop,while,for,in,struct,enum,impl,trait,pub,use,mod,self,super,crate,async,await,move,ref,and,or,not,true,false,None,Some,Ok,Err,where,as,const,static,type,unsafe,extern},
  keywordstyle=\color{blue}\bfseries,
  ndkeywords={bool,u8,u16,u32,u64,i32,i64,usize,isize,f64,String,Vec,Box,Arc,Mutex,Option,Result,Self},
  ndkeywordstyle=\color{purple}\bfseries,
  commentstyle=\color{gray}\itshape,
  stringstyle=\color{red},
  sensitive=true,
  morecomment=[l]{//},
  morecomment=[s]{/*}{*/},
  morestring=[b]",
}

\lstset{
  basicstyle=\ttfamily\scriptsize,
  keywordstyle=\color{blue},
  commentstyle=\color{gray},
  stringstyle=\color{red},
  breaklines=true,
  frame=single,
  columns=fullflexible,
  language=Rust
}

\title[StarryOS 大作业]{StarryOS 大作业总结\\[2mm] 支持 jcode 运行的内核调试与修复}
\author[姓名]{学生：姓名\\[5mm] 学号：XXXXXXXX}
\institute[清华大学]{\small 清华大学}
\date{\small \vskip -10pt 2026 年 6 月}

\begin{document}

% ============================================================
% 封面
% ============================================================
\begin{frame}
  \maketitle
\end{frame}

% ============================================================
% 目录
% ============================================================
\section*{目录}
\frame{
  \frametitle{\secname}
  \tableofcontents[hideallsubsections]
}

% ============================================================
% 1. 项目背景
% ============================================================
\section{项目背景}

\begin{frame}{jcode 是什么}
  \begin{block}{jcode：一个 TUI（终端用户界面）应用}
    \begin{itemize}
      \setlength{\itemsep}{4pt}
      \item 架构：TUI 进程 + serve 进程，通过 Unix socket 通信
      \item 技术栈：ratatui 框架、Tokio 异步运行时、reqwest HTTP 客户端
      \item 功能：AI 代码助手，支持鼠标点击、窗口调整、网络请求
    \end{itemize}
  \end{block}
  \begin{block}{对内核的要求}
    \begin{itemize}
      \setlength{\itemsep}{4pt}
      \item ext4 文件系统（读写文件）
      \item Unix domain socket（进程间通信）
      \item epoll + 非阻塞 I/O（Tokio 事件循环）
      \item TCP/IP 网络（HTTPS 请求）
      \item TTY 子系统（终端交互、鼠标、窗口大小）
    \end{itemize}
  \end{block}
  \centering
  \textbf{jcode 是一个"内核压力测试"——用到了几乎所有内核子系统}
\end{frame}

\begin{frame}{运行架构}
  \begin{tikzpicture}[
    node distance=0.8cm and 1.2cm,
    box/.style={rectangle, draw, rounded corners, minimum width=2.2cm, minimum height=0.8cm, align=center, font=\small},
    arr/.style={-{Stealth[length=3mm]}, thick}
  ]
    \node[box] (host) {宿主机终端};
    \node[box, below=of host] (ostool) {ostool\\{\tiny (用户态工具)}};
    \node[box, below=of ostool] (qemu) {QEMU 虚拟机};
    \node[box, below left=0.8cm and 0.5cm of qemu] (kernel) {StarryOS 内核};
    \node[box, below right=0.8cm and 0.5cm of qemu] (jcode) {jcode 进程};

    \draw[arr] (host) -- (ostool);
    \draw[arr] (ostool) -- (qemu);
    \draw[arr] (qemu) -- (kernel);
    \draw[arr] (qemu) -- (jcode);

    \node[font=\tiny, right=0.1cm of ostool] {键盘/鼠标};
    \node[font=\tiny, right=0.1cm of qemu] {串口 stdin/stdout};
  \end{tikzpicture}
\end{frame}

% ============================================================
% 2. 修复概览
% ============================================================
\section{修复概览}

\begin{frame}{修复概览：17 个 bug，7 个内核模块}
  \begin{table}[htbp]
    \scriptsize
    \centering
    \begin{tabular}{lcc}
      \toprule
      模块 & 修复数 & 典型问题 \\
      \midrule
      ext4 文件系统 & 3 & 块分配失败、rename 后写入、日志损坏 \\
      内存管理 & 1 & COW 缺页 panic \\
      epoll 事件通知 & 2 & ET 竞争窗口、NoEvent 死循环 \\
      任务调度 & 3 & device\_mask、抢占式唤醒、双重入队 \\
      网络协议栈 & 6 & ARP 队列、Unix socket accept/EOF \\
      TTY 终端 & 3 & 终端尺寸、SIGWINCH、批量输出 \\
      ostool & 2 & 鼠标事件转发、Moved 过滤 \\
      \bottomrule
    \end{tabular}
  \end{table}
\end{frame}

% ============================================================
% 3. ext4 文件系统
% ============================================================
\section{ext4 文件系统}

\begin{frame}[fragile]{Fix 1a：块分配跳过不一致块组}
  \begin{block}{问题}
    Alpine rootfs 非正常关机后，块组描述符（BGD）声称有空闲块，但位图全满。
    原始代码遇到 ENOSPC 直接报错 → 整个文件系统报"磁盘已满"。
  \end{block}
  \begin{block}{修复}
    对 ENOSPC 单独处理，跳过不一致的块组，继续尝试后续块组：
    \begin{lstlisting}[language=Rust]
let alloc = match alloc_res {
    Ok(a) => a,
    Err(e) if e.code == Errno::ENOSPC => {
        warn!("group={} inconsistent, skip", group_idx);
        continue;  // skip, try next group
    }
    Err(e) => return Err(e),
};\end{lstlisting}
  \end{block}
\end{frame}

\begin{frame}[fragile]{Fix 1b：rename 后写入失效}
  \begin{block}{问题}
    jcode 使用原子替换保存文件（write tmp $\to$ rename tmp $\to$ target）。
    VFS 的 Inode 缓存了 path，rename 后旧路径失效 $\to$ Drop 时 sync 用旧路径写 $\to$ ENOENT。
  \end{block}
  \begin{block}{修复}
    \begin{lstlisting}[language=Rust]
// old: use path (stale after rename)
rsext4::write_file(dev, fs, &self.path, offset, buf)
// new: use inode number (never changes)
rsext4::write_inode_data(dev, fs, self.ino, offset, buf)\end{lstlisting}
  \end{block}
  \begin{block}{根因}
    Inode 不应该缓存 path。path 是 DirEntry（目录树）的职责，不是 Inode（磁盘实体）的职责。
  \end{block}
\end{frame}

\begin{frame}{Fix 1c：JBD2 日志超级块损坏}
  \begin{block}{问题}
    Alpine rootfs 的 JBD2 超级块损坏：s\_maxlen=1（实际 8192），s\_start=8192（越界）。
    导致挂载时日志重放失败，文件系统无法挂载。
  \end{block}
  \begin{block}{修复}
    \begin{itemize}
      \item s\_maxlen $<$ s\_first（明显损坏）$\to$ 用物理分区实际长度
      \item s\_start $\ge$ journal\_blocks.len()（越界）$\to$ 视为"日志干净"，跳过重放
    \end{itemize}
  \end{block}
\end{frame}

% ============================================================
% 4. 内存管理
% ============================================================
\section{内存管理与 COW}

\begin{frame}{Fix 2：内核态 COW 缺页 panic}
  \begin{block}{问题}
    \texttt{get\_as\_mut()} 获取用户空间 \texttt{\&mut} 引用。
    另一个线程调用 \texttt{fork()} 将该页标记为 COW 只读。
    内核写入 $\to$ CPU \#PF（写保护异常）$\to$ panic。
  \end{block}
  \begin{block}{修复（两层）}
    \begin{enumerate}
      \setlength{\itemsep}{4pt}
      \item \textbf{主修复}：epoll/poll 改用内核缓冲区 + \texttt{vm\_write\_slice} 安全写回
      \item \textbf{兜底}：扩展 \texttt{handle\_page\_fault}，内核态缺页时也能处理用户空间 COW
    \end{enumerate}
  \end{block}
  \begin{block}{思考}
    \texttt{get\_as\_mut()} 是不安全的访问方式，应该逐步替换为 \texttt{vm\_read/write\_slice}。
  \end{block}
\end{frame}

% ============================================================
% 5. epoll 事件通知
% ============================================================
\section{epoll 事件通知}

\begin{frame}{Fix 3a：EPOLLET 竞争窗口}
  \begin{block}{问题}
    ET 模式要求"每次状态变化恰好通知一次"。
    \texttt{mark\_not\_in\_queue()} 和 \texttt{register\_waker\_only()} 之间数据到达 $\to$ 通知丢失。
  \end{block}
  \begin{block}{修复}
    注册新 Waker 后再次检查文件状态。如果数据已在缓冲区，直接将 interest 入队并唤醒 epoll\_wait。
  \end{block}
  \begin{block}{Fix 3b：NoEvent 路径死循环}
    TCP socket 的 EPOLLOUT 始终就绪。虚假唤醒 $\to$ \texttt{check\_and\_register\_waker} $\to$ wake $\to$ re-queue $\to$ 死循环。
    修复：NoEvent 路径只注册 Waker，不主动 poll/唤醒。
  \end{block}
\end{frame}

% ============================================================
% 6. 异步任务调度
% ============================================================
\section{异步任务调度}

\begin{frame}{Fix 4：device\_mask 与 Waker 注册}
  \begin{block}{核心概念}
    TCP socket 的 \texttt{device\_mask} 决定 Waker 注册到哪个网卡设备的 PollSet 上。
    默认为 0，只有在 \texttt{start\_connect()} 内部才被设置。
  \end{block}
  \begin{block}{问题}
    Tokio 的典型流程：\texttt{epoll\_ctl(ADD)} $\to$ \texttt{connect()}。
    epoll\_ctl 时 device\_mask=0，InterestWaker 注册到空 mask $\to$ 等于没注册。
    connect 设好 mask 后，原来注册的 InterestWaker 不会自动迁移。
  \end{block}
  \begin{block}{修复}
    poll\_io 无论阻塞/非阻塞都先注册 Waker。connect 在 mask 设好后再注册一次。
  \end{block}
\end{frame}

\begin{frame}{Fix 5a/5b：抢占式唤醒与双重入队}
  \begin{block}{Fix 5a：抢占式唤醒}
    原始 \texttt{unblock\_task(task, false)} 不设置抢占标志，键盘响应延迟 20ms。
    修复：改为 \texttt{unblock\_task(task, true)} 立即抢占。
  \end{block}
  \begin{block}{Fix 5b：双重入队}
    抢占后任务在 \texttt{wait\_until} 中被唤醒，条件仍不满足（锁未释放），再次入队 $\to$ 重复条目。
    修复：入队前检查 \texttt{in\_wait\_queue} 标志。
  \end{block}
  \begin{block}{为什么条件仍不满足？}
    \texttt{notify\_one(resched=true)} 在持有锁的代码中调用。
    唤醒 T 后锁还没释放。抢占式唤醒让 T 立刻运行，但锁还在别人手里。
  \end{block}
\end{frame}

% ============================================================
% 7. 网络协议栈
% ============================================================
\section{网络协议栈}

\begin{frame}{ARP 三个修复}
  \begin{block}{Fix 6a：队列扩容 32 $\to$ 128}
    jcode 启动时并发建立 10-20 个 TCP 连接，32 个槽位耗尽，SYN 包被丢弃。
  \end{block}
  \begin{block}{Fix 6b：TTL 延长 60s $\to$ 300s}
    AI 响应可能超过 60s，ARP 缓存过期后所有 ACK 包重新排队，压爆队列。
  \end{block}
  \begin{block}{Fix 6c：全队列扫描}
    原始代码只处理队头。队列 [A, B, A, A]（A 已解析，B 未解析）$\to$ 只发第一个 A，后面的 A 永远卡在 B 后面。
    修复：扫描整个队列，发送所有已解析的包。
  \end{block}
\end{frame}

\begin{frame}{Unix socket 两个修复}
  \begin{block}{Fix 6d：非阻塞 accept}
    \texttt{accept()} 内部用 \texttt{rx.recv().await}，完全忽略 O\_NONBLOCK。
    Tokio 用 EPOLLET，收到 EPOLLIN 后循环 accept 直到 EAGAIN，但旧代码永远不返回 $\to$ Tokio 冻结。
    修复：添加 \texttt{try\_accept()}，无连接时立即返回 WouldBlock。
  \end{block}
  \begin{block}{Fix 6e：EOF 信号顺序}
    Drop 时 HeapProd 还活着（Rust 字段 drop 在函数体之后）。
    对端 poll 看到 \texttt{write\_is\_held()=true} + 空缓冲区 $\to$ 误判为"没关闭"。
    修复：先设 \texttt{my\_tx\_closed} 标志再 wake。
  \end{block}
\end{frame}

% ============================================================
% 8. TTY 终端
% ============================================================
\section{TTY 终端}

\begin{frame}{Fix 7a/7b/7c：终端子系统三个修复}
  \begin{block}{Fix 7a：查询真实终端尺寸}
    ANSI CPR 技术：\texttt{\textbackslash x1b[9999;9999H} 移动光标到边界 $\to$
    \texttt{\textbackslash x1b[6n} 请求位置 $\to$ 终端回复 \texttt{\textbackslash x1b[rows;colsR}。
    默认值改为 24x80（VT100 标准）。
  \end{block}
  \begin{block}{Fix 7b：SIGWINCH 信号}
    \texttt{TIOCSWINSZ} 更新窗口大小后，检查尺寸变化，给前台进程组发 SIGWINCH。
    TUI 收到信号 $\to$ 重新读取尺寸 $\to$ 刷新界面。
  \end{block}
  \begin{block}{Fix 7c：批量 UART 输出}
    每个 \texttt{\textbackslash n} 两次写操作 $\to$ 锁获取/释放 $\to$ 逐行渲染 $\to$ 闪烁。
    修复：先收集到 Vec，一次写入，一帧原子发送。
  \end{block}
\end{frame}

% ============================================================
% 9. ostool 鼠标转发
% ============================================================
\section{ostool 鼠标转发}

\begin{frame}{Fix 8a/8b：ostool 鼠标事件}
  \begin{block}{ostool 是什么}
    宿主机上的用户态工具，不在 QEMU 内。读取终端键盘/鼠标事件 $\to$ 写入 QEMU stdin $\to$
    读取 QEMU stdout $\to$ 显示在宿主机终端。
  \end{block}
  \begin{block}{Fix 8a：鼠标事件转发}
    问题：ostool 静默丢弃鼠标事件。
    修复：编码为 SGR 格式（\texttt{\textbackslash x1b[<Cb;Cx;CyM/m}）转发给 guest。
  \end{block}
  \begin{block}{Fix 8b：过滤 Moved 事件}
    每像素一个 Moved 事件 $\to$ TUI 每帧重绘 $\to$ UART 饱和 $\to$ 冻结。
    修复：丢弃 Moved，只转发点击/滚动/拖拽。
  \end{block}
\end{frame}

% ============================================================
% 10. 关键收获
% ============================================================
\section{关键收获}

\begin{frame}{关键收获}
  \begin{enumerate}
    \setlength{\itemsep}{6pt}
    \item \textbf{VFS 分层设计}：DirEntry 树（内存缓存）vs Inode（磁盘实体）vs ext4 目录块。Inode 不应该缓存 path。
    \item \textbf{异步执行模型}：Rust Future = poll 状态机，block\_on = 驱动器，poll\_io = 通用 I/O 等待。
    \item \textbf{Waker 注册的生命周期}：device\_mask 在 connect 时才设置，epoll\_ctl(ADD) 可能早于 connect $\to$ 注册无效。
    \item \textbf{抢占式调度的副作用}：提高响应速度，但引入竞争条件——任务可能在条件还不满足时就被唤醒。
    \item \textbf{Rust Drop 顺序}：函数体先执行，字段后 drop $\to$ HeapProd 在 wake 时还活着，需要原子标志绕过。
  \end{enumerate}
\end{frame}

% ============================================================
% 11. AI 时代 OS 教学思考
% ============================================================
\section{AI 时代 OS 教学思考}

\begin{frame}{AI 时代 OS 教学思考}
  \begin{block}{传统实验的挑战}
    AI 能自动完成"实现 syscall"类任务。单模块实验价值下降。
  \end{block}
  \begin{block}{建议方向}
    \begin{enumerate}
      \setlength{\itemsep}{4pt}
      \item \textbf{从"实现"转向"调试"}：给有 bug 的内核，让学生定位问题
      \item \textbf{跨模块系统性问题}：涉及多个子系统的交互 bug
      \item \textbf{真实应用驱动}：用 jcode 这类真实应用触发边界条件
      \item \textbf{性能调优}：从"功能正确"到"性能优化"
      \item \textbf{安全性审查}：学会审查 AI 生成的代码
    \end{enumerate}
  \end{block}
\end{frame}

% ============================================================
% 致谢
% ============================================================
\begin{frame}
  \begin{center}
    {\Huge\calligra Thanks for your attention!}
    \vspace{1cm}

    {\Huge Q \& A}
  \end{center}
\end{frame}

\end{document}
